%%%%%%%%%%%%%%%%%%%%%%%%%%%%%%%%%%%%%%%%%%%%%%%%%%%%%%%%%%%%
% 6.3 RECURRENCE RELATIONS
% The Composition of Advanced Mathematics
%%%%%%%%%%%%%%%%%%%%%%%%%%%%%%%%%%%%%%%%%%%%%%%%%%%%%%%%%%%%

\section{Recurrence Relations}
\label{sec:recurrence-relations}

\prerequisite{
Sequences, functions, algebraic manipulation, induction, basic counting,
and the counting principles developed in
Section~\ref{sec:basic-counting}.
}

\learningobjectives{
By the end of this section, the reader should be able to:
\begin{itemize}
    \item define and interpret recurrence relations;
    \item identify initial conditions and determine whether they specify
          a sequence uniquely;
    \item distinguish linear and nonlinear recurrences;
    \item distinguish homogeneous and nonhomogeneous linear recurrences;
    \item classify recurrences by order and coefficient type;
    \item compute sequence terms recursively;
    \item derive recurrence relations from combinatorial problems;
    \item solve first-order linear recurrences;
    \item solve linear homogeneous recurrences with constant coefficients;
    \item use characteristic equations and characteristic roots;
    \item handle repeated characteristic roots;
    \item solve common nonhomogeneous recurrences;
    \item apply substitution and iteration methods;
    \item recognize divide-and-conquer recurrences;
    \item apply the Master Theorem in appropriate cases;
    \item use recursion trees to analyze recurrence growth;
    \item recognize Fibonacci-type and Catalan-type recurrences;
    \item understand recurrences arising from dynamic programming;
    \item distinguish exact solutions from asymptotic estimates;
    \item verify proposed recurrence solutions.
\end{itemize}
}

%%%%%%%%%%%%%%%%%%%%%%%%%%%%%%%%%%%%%%%%%%%%%%%%%%%%%%%%%%%%
\subsection{What Is a Recurrence Relation?}
\label{subsec:what-is-recurrence}
%%%%%%%%%%%%%%%%%%%%%%%%%%%%%%%%%%%%%%%%%%%%%%%%%%%%%%%%%%%%

A recurrence relation defines terms of a sequence using earlier terms.

\begin{definitionbox}{Recurrence Relation}
A recurrence relation is an equation that expresses one or more terms
of a sequence in terms of preceding terms.
\end{definitionbox}

For example,

\[
a_n=a_{n-1}+3
\]

defines each term from the previous one.

If

\[
a_0=2,
\]

then

\[
a_1=5,
\qquad
a_2=8,
\qquad
a_3=11,
\]

and so on.

The recurrence and its initial condition together determine the
sequence.

%%%%%%%%%%%%%%%%%%%%%%%%%%%%%%%%%%%%%%%%%%%%%%%%%%%%%%%%%%%%
\subsection{Initial Conditions}
\label{subsec:initial-conditions}
%%%%%%%%%%%%%%%%%%%%%%%%%%%%%%%%%%%%%%%%%%%%%%%%%%%%%%%%%%%%

A recurrence relation alone often does not determine a unique sequence.

For example,

\[
a_n=a_{n-1}+3
\]

could generate infinitely many different sequences depending on the
value of \(a_0\).

An initial condition specifies one or more starting values.

For a first-order recurrence, one initial value is often sufficient.

For a second-order recurrence such as

\[
a_n=a_{n-1}+a_{n-2},
\]

two initial values are generally required.

%%%%%%%%%%%%%%%%%%%%%%%%%%%%%%%%%%%%%%%%%%%%%%%%%%%%%%%%%%%%
\subsection{Order of a Recurrence}
\label{subsec:order-recurrence}
%%%%%%%%%%%%%%%%%%%%%%%%%%%%%%%%%%%%%%%%%%%%%%%%%%%%%%%%%%%%

\begin{definitionbox}{Order of a Recurrence}
The order of a recurrence is the largest number of steps backward
required to determine the current term.
\end{definitionbox}

For example,

\[
a_n=2a_{n-1}+1
\]

is first order.

The recurrence

\[
a_n=3a_{n-1}-2a_{n-2}
\]

is second order.

The recurrence

\[
a_n=a_{n-1}+a_{n-3}
\]

is third order.

%%%%%%%%%%%%%%%%%%%%%%%%%%%%%%%%%%%%%%%%%%%%%%%%%%%%%%%%%%%%
\subsection{Recursive Computation}
\label{subsec:recursive-computation}
%%%%%%%%%%%%%%%%%%%%%%%%%%%%%%%%%%%%%%%%%%%%%%%%%%%%%%%%%%%%

Suppose

\[
a_0=1
\]

and

\[
a_n=2a_{n-1}+1.
\]

Then

\[
a_1=2(1)+1=3,
\]

\[
a_2=2(3)+1=7,
\]

\[
a_3=2(7)+1=15.
\]

The recurrence gives an algorithm for computing terms one at a time.

However, a closed formula may allow direct computation of \(a_n\)
without first calculating every preceding term.

%%%%%%%%%%%%%%%%%%%%%%%%%%%%%%%%%%%%%%%%%%%%%%%%%%%%%%%%%%%%
\subsection{Explicit and Recursive Descriptions}
\label{subsec:explicit-recursive}
%%%%%%%%%%%%%%%%%%%%%%%%%%%%%%%%%%%%%%%%%%%%%%%%%%%%%%%%%%%%

A sequence may be described recursively or explicitly.

For example,

\[
a_n=a_{n-1}+3,
\qquad
a_0=2
\]

is recursive.

The equivalent explicit formula is

\[
\boxed{
a_n=2+3n.
}
\]

The recursive form emphasizes how the sequence evolves.

The explicit form emphasizes the value at a particular index.

%%%%%%%%%%%%%%%%%%%%%%%%%%%%%%%%%%%%%%%%%%%%%%%%%%%%%%%%%%%%
\subsection{Worked Example: Computing Terms}
\label{subsec:worked-computing-recurrence}
%%%%%%%%%%%%%%%%%%%%%%%%%%%%%%%%%%%%%%%%%%%%%%%%%%%%%%%%%%%%

\begin{workedexamplebox}{Evaluating a Recurrence}

Suppose

\[
a_0=4
\]

and

\[
a_n=3a_{n-1}-2.
\]

Then

\[
a_1=3(4)-2=10,
\]

\[
a_2=3(10)-2=28,
\]

\[
a_3=3(28)-2=82.
\]

\begin{answerbox}
\[
\boxed{
a_1=10,\qquad
a_2=28,\qquad
a_3=82.
}
\]
\end{answerbox}

\end{workedexamplebox}

%%%%%%%%%%%%%%%%%%%%%%%%%%%%%%%%%%%%%%%%%%%%%%%%%%%%%%%%%%%%
\subsection{Linear Recurrences}
\label{subsec:linear-recurrences}
%%%%%%%%%%%%%%%%%%%%%%%%%%%%%%%%%%%%%%%%%%%%%%%%%%%%%%%%%%%%

A recurrence is linear when the sequence terms appear only to the first
power and are not multiplied together.

A general linear recurrence of order \(k\) has the form

\[
\boxed{
a_n
=
c_1(n)a_{n-1}
+
c_2(n)a_{n-2}
+\cdots+
c_k(n)a_{n-k}
+
f(n).
}
\]

The functions

\[
c_1(n),\ldots,c_k(n)
\]

are coefficients.

The function

\[
f(n)
\]

is the forcing term.

%%%%%%%%%%%%%%%%%%%%%%%%%%%%%%%%%%%%%%%%%%%%%%%%%%%%%%%%%%%%
\subsection{Linear Homogeneous Recurrences}
\label{subsec:linear-homogeneous-recurrences}
%%%%%%%%%%%%%%%%%%%%%%%%%%%%%%%%%%%%%%%%%%%%%%%%%%%%%%%%%%%%

A linear recurrence is homogeneous when the forcing term is zero.

Thus,

\[
\boxed{
a_n
=
c_1a_{n-1}
+
c_2a_{n-2}
+\cdots+
c_ka_{n-k}
}
\]

is homogeneous when the coefficients are constants.

For example,

\[
a_n=5a_{n-1}-6a_{n-2}
\]

is a second-order linear homogeneous recurrence with constant
coefficients.

%%%%%%%%%%%%%%%%%%%%%%%%%%%%%%%%%%%%%%%%%%%%%%%%%%%%%%%%%%%%
\subsection{Linear Nonhomogeneous Recurrences}
\label{subsec:linear-nonhomogeneous-recurrences}
%%%%%%%%%%%%%%%%%%%%%%%%%%%%%%%%%%%%%%%%%%%%%%%%%%%%%%%%%%%%

A recurrence such as

\[
\boxed{
a_n=2a_{n-1}+3
}
\]

is linear but nonhomogeneous because of the nonzero forcing term.

Similarly,

\[
a_n
=
4a_{n-1}
-
4a_{n-2}
+
n^2
\]

is nonhomogeneous.

The solution generally consists of

\[
\boxed{
\text{homogeneous solution}
+
\text{particular solution}.
}
\]

%%%%%%%%%%%%%%%%%%%%%%%%%%%%%%%%%%%%%%%%%%%%%%%%%%%%%%%%%%%%
\subsection{Nonlinear Recurrences}
\label{subsec:nonlinear-recurrences}
%%%%%%%%%%%%%%%%%%%%%%%%%%%%%%%%%%%%%%%%%%%%%%%%%%%%%%%%%%%%

Recurrences may also be nonlinear.

Examples include

\[
a_n=(a_{n-1})^2+1,
\]

\[
a_n=a_{n-1}a_{n-2},
\]

and

\[
a_n=\sqrt{1+a_{n-1}}.
\]

Nonlinear recurrences usually require methods different from
characteristic equations.

%%%%%%%%%%%%%%%%%%%%%%%%%%%%%%%%%%%%%%%%%%%%%%%%%%%%%%%%%%%%
\subsection{First-Order Arithmetic Recurrences}
\label{subsec:first-order-arithmetic}
%%%%%%%%%%%%%%%%%%%%%%%%%%%%%%%%%%%%%%%%%%%%%%%%%%%%%%%%%%%%

Consider

\[
a_n=a_{n-1}+d.
\]

Repeated substitution gives

\[
a_n=a_{n-2}+2d,
\]

then

\[
a_n=a_{n-3}+3d.
\]

Continuing,

\[
a_n=a_0+nd.
\]

Thus,

\[
\boxed{
a_n=a_0+nd.
}
\]

This is the familiar arithmetic sequence.

%%%%%%%%%%%%%%%%%%%%%%%%%%%%%%%%%%%%%%%%%%%%%%%%%%%%%%%%%%%%
\subsection{First-Order Geometric Recurrences}
\label{subsec:first-order-geometric}
%%%%%%%%%%%%%%%%%%%%%%%%%%%%%%%%%%%%%%%%%%%%%%%%%%%%%%%%%%%%

Consider

\[
a_n=ra_{n-1}.
\]

Repeated substitution gives

\[
a_n=r^2a_{n-2},
\]

then

\[
a_n=r^3a_{n-3}.
\]

Continuing,

\[
\boxed{
a_n=a_0r^n.
}
\]

This is the geometric sequence.

%%%%%%%%%%%%%%%%%%%%%%%%%%%%%%%%%%%%%%%%%%%%%%%%%%%%%%%%%%%%
\subsection{Worked Example: Geometric Recurrence}
\label{subsec:worked-geometric-recurrence}
%%%%%%%%%%%%%%%%%%%%%%%%%%%%%%%%%%%%%%%%%%%%%%%%%%%%%%%%%%%%

\begin{workedexamplebox}{Solving a Multiplicative Recurrence}

Suppose

\[
a_0=5
\]

and

\[
a_n=3a_{n-1}.
\]

The recurrence is geometric with ratio \(3\).

Therefore,

\begin{answerbox}
\[
\boxed{
a_n=5\cdot3^n.
}
\]
\end{answerbox}

\end{workedexamplebox}

%%%%%%%%%%%%%%%%%%%%%%%%%%%%%%%%%%%%%%%%%%%%%%%%%%%%%%%%%%%%
\subsection{First-Order Linear Recurrences}
\label{subsec:first-order-linear}
%%%%%%%%%%%%%%%%%%%%%%%%%%%%%%%%%%%%%%%%%%%%%%%%%%%%%%%%%%%%

Consider

\[
\boxed{
a_n=ra_{n-1}+b,
}
\]

where \(r\) and \(b\) are constants.

When

\[
r\neq1,
\]

the recurrence has equilibrium value

\[
L=\frac{b}{1-r}.
\]

Subtract \(L\) from both sides:

\[
a_n-L
=
r(a_{n-1}-L).
\]

Therefore,

\[
a_n-L
=
r^n(a_0-L).
\]

Hence,

\[
\boxed{
a_n
=
L+(a_0-L)r^n,
\qquad
L=\frac{b}{1-r}.
}
\]

%%%%%%%%%%%%%%%%%%%%%%%%%%%%%%%%%%%%%%%%%%%%%%%%%%%%%%%%%%%%
\subsection{Worked Example: First-Order Linear Recurrence}
\label{subsec:worked-first-order-linear}
%%%%%%%%%%%%%%%%%%%%%%%%%%%%%%%%%%%%%%%%%%%%%%%%%%%%%%%%%%%%

\begin{workedexamplebox}{Solving \(a_n=2a_{n-1}+3\)}

Suppose

\[
a_0=1
\]

and

\[
a_n=2a_{n-1}+3.
\]

The equilibrium satisfies

\[
L=2L+3,
\]

so

\[
L=-3.
\]

Then

\[
a_n+3
=
2(a_{n-1}+3).
\]

Therefore,

\[
a_n+3
=
2^n(a_0+3).
\]

Since \(a_0=1\),

\[
a_n+3=4\cdot2^n.
\]

\begin{answerbox}
\[
\boxed{
a_n=4\cdot2^n-3.
}
\]
\end{answerbox}

\end{workedexamplebox}

%%%%%%%%%%%%%%%%%%%%%%%%%%%%%%%%%%%%%%%%%%%%%%%%%%%%%%%%%%%%
\subsection{Iteration Method}
\label{subsec:iteration-method}
%%%%%%%%%%%%%%%%%%%%%%%%%%%%%%%%%%%%%%%%%%%%%%%%%%%%%%%%%%%%

A recurrence can often be solved by repeatedly substituting earlier
expressions.

Suppose

\[
a_n=2a_{n-1}+1.
\]

Then

\[
a_n
=
2(2a_{n-2}+1)+1
=
2^2a_{n-2}+2+1.
\]

Repeating,

\[
a_n
=
2^ka_{n-k}
+
\sum_{j=0}^{k-1}2^j.
\]

Taking \(k=n\),

\[
a_n
=
2^na_0
+
\sum_{j=0}^{n-1}2^j.
\]

Using the geometric sum,

\[
\sum_{j=0}^{n-1}2^j
=
2^n-1.
\]

Hence,

\[
\boxed{
a_n
=
2^na_0+2^n-1.
}
\]

%%%%%%%%%%%%%%%%%%%%%%%%%%%%%%%%%%%%%%%%%%%%%%%%%%%%%%%%%%%%
\subsection{Second-Order Homogeneous Recurrences}
\label{subsec:second-order-homogeneous}
%%%%%%%%%%%%%%%%%%%%%%%%%%%%%%%%%%%%%%%%%%%%%%%%%%%%%%%%%%%%

Consider

\[
\boxed{
a_n
=
c_1a_{n-1}
+
c_2a_{n-2}.
}
\]

We seek solutions of the form

\[
a_n=r^n.
\]

Substitution gives

\[
r^n
=
c_1r^{n-1}
+
c_2r^{n-2}.
\]

Assuming \(r\neq0\), divide by \(r^{n-2}\):

\[
r^2
=
c_1r+c_2.
\]

Thus,

\[
\boxed{
r^2-c_1r-c_2=0.
}
\]

This polynomial equation is called the characteristic equation.

%%%%%%%%%%%%%%%%%%%%%%%%%%%%%%%%%%%%%%%%%%%%%%%%%%%%%%%%%%%%
\subsection{Characteristic Roots}
\label{subsec:characteristic-roots}
%%%%%%%%%%%%%%%%%%%%%%%%%%%%%%%%%%%%%%%%%%%%%%%%%%%%%%%%%%%%

The solutions of the characteristic equation are called characteristic
roots.

Suppose the roots are distinct:

\[
r_1\neq r_2.
\]

Then the general solution has the form

\[
\boxed{
a_n
=
A r_1^n
+
B r_2^n.
}
\]

The constants \(A\) and \(B\) are determined by the initial conditions.

%%%%%%%%%%%%%%%%%%%%%%%%%%%%%%%%%%%%%%%%%%%%%%%%%%%%%%%%%%%%
\subsection{Worked Example: Distinct Characteristic Roots}
\label{subsec:worked-distinct-roots}
%%%%%%%%%%%%%%%%%%%%%%%%%%%%%%%%%%%%%%%%%%%%%%%%%%%%%%%%%%%%

\begin{workedexamplebox}{Solving a Second-Order Recurrence}

Suppose

\[
a_n
=
5a_{n-1}
-
6a_{n-2},
\]

with

\[
a_0=1,
\qquad
a_1=4.
\]

The characteristic equation is

\[
r^2-5r+6=0.
\]

Factor:

\[
(r-2)(r-3)=0.
\]

Thus,

\[
r_1=2,
\qquad
r_2=3.
\]

The general solution is

\[
a_n=A2^n+B3^n.
\]

Use \(a_0=1\):

\[
A+B=1.
\]

Use \(a_1=4\):

\[
2A+3B=4.
\]

Solving gives

\[
A=-1,
\qquad
B=2.
\]

\begin{answerbox}
\[
\boxed{
a_n
=
-2^n+2\cdot3^n.
}
\]
\end{answerbox}

\end{workedexamplebox}

%%%%%%%%%%%%%%%%%%%%%%%%%%%%%%%%%%%%%%%%%%%%%%%%%%%%%%%%%%%%
\subsection{Repeated Characteristic Roots}
\label{subsec:repeated-characteristic-roots}
%%%%%%%%%%%%%%%%%%%%%%%%%%%%%%%%%%%%%%%%%%%%%%%%%%%%%%%%%%%%

Suppose the characteristic equation has a repeated root \(r\).

For a second-order recurrence, the usual form

\[
A r^n+B r^n
\]

does not produce two independent solutions.

Instead, the general solution is

\[
\boxed{
a_n
=
(A+Bn)r^n.
}
\]

More generally, if \(r\) has multiplicity \(m\), the associated terms
are

\[
\boxed{
r^n,
\quad
nr^n,
\quad
n^2r^n,
\quad
\ldots,
\quad
n^{m-1}r^n.
}
\]

%%%%%%%%%%%%%%%%%%%%%%%%%%%%%%%%%%%%%%%%%%%%%%%%%%%%%%%%%%%%
\subsection{Worked Example: Repeated Root}
\label{subsec:worked-repeated-root}
%%%%%%%%%%%%%%%%%%%%%%%%%%%%%%%%%%%%%%%%%%%%%%%%%%%%%%%%%%%%

\begin{workedexamplebox}{A Repeated Characteristic Root}

Suppose

\[
a_n=6a_{n-1}-9a_{n-2},
\]

with

\[
a_0=2,
\qquad
a_1=9.
\]

The characteristic equation is

\[
r^2-6r+9=0.
\]

Thus,

\[
(r-3)^2=0.
\]

The repeated root is

\[
r=3.
\]

Therefore,

\[
a_n=(A+Bn)3^n.
\]

Using \(a_0=2\),

\[
A=2.
\]

Using \(a_1=9\),

\[
(2+B)3=9,
\]

so

\[
B=1.
\]

\begin{answerbox}
\[
\boxed{
a_n=(2+n)3^n.
}
\]
\end{answerbox}

\end{workedexamplebox}

%%%%%%%%%%%%%%%%%%%%%%%%%%%%%%%%%%%%%%%%%%%%%%%%%%%%%%%%%%%%
\subsection{Higher-Order Linear Recurrences}
\label{subsec:higher-order-linear-recurrences}
%%%%%%%%%%%%%%%%%%%%%%%%%%%%%%%%%%%%%%%%%%%%%%%%%%%%%%%%%%%%

Consider

\[
a_n
=
c_1a_{n-1}
+
c_2a_{n-2}
+\cdots+
c_ka_{n-k}.
\]

Seeking solutions

\[
a_n=r^n
\]

produces the characteristic equation

\[
\boxed{
r^k
-
c_1r^{k-1}
-
c_2r^{k-2}
-\cdots-
c_k
=
0.
}
\]

Each root contributes terms according to its multiplicity.

The initial conditions determine the unknown coefficients.

%%%%%%%%%%%%%%%%%%%%%%%%%%%%%%%%%%%%%%%%%%%%%%%%%%%%%%%%%%%%
\subsection{Complex Characteristic Roots}
\label{subsec:complex-characteristic-roots}
%%%%%%%%%%%%%%%%%%%%%%%%%%%%%%%%%%%%%%%%%%%%%%%%%%%%%%%%%%%%

Characteristic equations may have complex roots.

Suppose the roots are

\[
r=\rho e^{i\theta}
\]

and

\[
\overline{r}
=
\rho e^{-i\theta}.
\]

For a recurrence with real coefficients, the resulting real solution
can be written

\[
\boxed{
a_n
=
\rho^n
\left(
A\cos(n\theta)
+
B\sin(n\theta)
\right).
}
\]

Here,

\[
\rho
\]

is the Greek letter ``rho'' and represents the magnitude of the root.

The angle

\[
\theta
\]

is the Greek letter ``theta.''

%%%%%%%%%%%%%%%%%%%%%%%%%%%%%%%%%%%%%%%%%%%%%%%%%%%%%%%%%%%%
\subsection{The Fibonacci Recurrence}
\label{subsec:fibonacci-recurrence}
%%%%%%%%%%%%%%%%%%%%%%%%%%%%%%%%%%%%%%%%%%%%%%%%%%%%%%%%%%%%

The Fibonacci sequence is defined by

\[
F_0=0,
\qquad
F_1=1,
\]

and

\[
\boxed{
F_n
=
F_{n-1}
+
F_{n-2}.
}
\]

Its first terms are

\[
0,1,1,2,3,5,8,13,21,\ldots.
\]

The characteristic equation is

\[
r^2-r-1=0.
\]

The roots are

\[
\boxed{
\phi
=
\frac{1+\sqrt5}{2}
}
\]

and

\[
\boxed{
\psi
=
\frac{1-\sqrt5}{2}.
}
\]

The Greek letters are called ``phi'' and ``psi.''

%%%%%%%%%%%%%%%%%%%%%%%%%%%%%%%%%%%%%%%%%%%%%%%%%%%%%%%%%%%%
\subsection{Binet's Formula}
\label{subsec:binets-formula}
%%%%%%%%%%%%%%%%%%%%%%%%%%%%%%%%%%%%%%%%%%%%%%%%%%%%%%%%%%%%

Using the initial conditions, the Fibonacci sequence has the explicit
formula

\[
\boxed{
F_n
=
\frac{\phi^n-\psi^n}{\sqrt5}.
}
\]

This is called Binet's formula.

Because

\[
|\psi|<1,
\]

the term

\[
\psi^n
\]

approaches zero as \(n\) grows.

Therefore,

\[
\boxed{
F_n
\sim
\frac{\phi^n}{\sqrt5}.
}
\]

Thus Fibonacci numbers grow approximately exponentially with base
\(\phi\).

%%%%%%%%%%%%%%%%%%%%%%%%%%%%%%%%%%%%%%%%%%%%%%%%%%%%%%%%%%%%
\subsection{Deriving Recurrences from Counting Problems}
\label{subsec:deriving-recurrences-counting}
%%%%%%%%%%%%%%%%%%%%%%%%%%%%%%%%%%%%%%%%%%%%%%%%%%%%%%%%%%%%

A recurrence often appears when a combinatorial class can be divided
according to the first or last structural choice.

The general strategy is:

\begin{enumerate}
    \item identify a family of objects of size \(n\);
    \item split the objects into disjoint cases;
    \item reduce each case to smaller instances;
    \item translate the decomposition into an equation.
\end{enumerate}

This technique is one of the most important uses of recurrences in
combinatorics.

%%%%%%%%%%%%%%%%%%%%%%%%%%%%%%%%%%%%%%%%%%%%%%%%%%%%%%%%%%%%
\subsection{Worked Example: Binary Strings with No Consecutive Ones}
\label{subsec:binary-no-consecutive-ones}
%%%%%%%%%%%%%%%%%%%%%%%%%%%%%%%%%%%%%%%%%%%%%%%%%%%%%%%%%%%%

\begin{workedexamplebox}{A Fibonacci-Type Counting Problem}

Let \(a_n\) be the number of binary strings of length \(n\) containing
no two consecutive \(1\)'s.

Separate valid strings according to their ending.

If a valid string ends in \(0\), the preceding \(n-1\) positions form
any valid string of length \(n-1\).

This contributes

\[
a_{n-1}.
\]

If a valid string ends in \(1\), the previous position must be \(0\).

Removing the final \(01\) leaves any valid string of length \(n-2\).

This contributes

\[
a_{n-2}.
\]

Therefore,

\[
a_n=a_{n-1}+a_{n-2}.
\]

The initial values are

\[
a_0=1,
\qquad
a_1=2.
\]

\begin{answerbox}
\[
\boxed{
a_n=a_{n-1}+a_{n-2},
\qquad
a_0=1,
\qquad
a_1=2.
}
\]
\end{answerbox}

\end{workedexamplebox}

%%%%%%%%%%%%%%%%%%%%%%%%%%%%%%%%%%%%%%%%%%%%%%%%%%%%%%%%%%%%
\subsection{Tiling Recurrences}
\label{subsec:tiling-recurrences}
%%%%%%%%%%%%%%%%%%%%%%%%%%%%%%%%%%%%%%%%%%%%%%%%%%%%%%%%%%%%

Suppose a \(1\times n\) board is tiled using:

\[
1\times1
\]

squares and

\[
1\times2
\]

dominoes.

Let \(T_n\) denote the number of tilings.

A tiling ends either with:

\begin{itemize}
    \item one square, leaving a board of length \(n-1\); or
    \item one domino, leaving a board of length \(n-2\).
\end{itemize}

Thus,

\[
\boxed{
T_n=T_{n-1}+T_{n-2}.
}
\]

Again, a Fibonacci-type recurrence appears.

%%%%%%%%%%%%%%%%%%%%%%%%%%%%%%%%%%%%%%%%%%%%%%%%%%%%%%%%%%%%
\subsection{Recurrences from Restricted Strings}
\label{subsec:restricted-string-recurrences}
%%%%%%%%%%%%%%%%%%%%%%%%%%%%%%%%%%%%%%%%%%%%%%%%%%%%%%%%%%%%

Suppose strings are constructed under local restrictions.

The final symbol often determines which smaller string remains.

This makes recurrence relations particularly effective for:

\begin{itemize}
    \item forbidden substrings;
    \item run restrictions;
    \item adjacency conditions;
    \item color restrictions;
    \item tiling problems;
    \item finite-state models.
\end{itemize}

More complicated restrictions may require several related recurrence
sequences rather than one.

%%%%%%%%%%%%%%%%%%%%%%%%%%%%%%%%%%%%%%%%%%%%%%%%%%%%%%%%%%%%
\subsection{Systems of Recurrences}
\label{subsec:systems-recurrences}
%%%%%%%%%%%%%%%%%%%%%%%%%%%%%%%%%%%%%%%%%%%%%%%%%%%%%%%%%%%%

Sometimes one sequence is insufficient to describe the state of a
problem.

For example, let:

\[
a_n
=
\text{number of valid strings of length }n\text{ ending in }0,
\]

and

\[
b_n
=
\text{number ending in }1.
\]

One may obtain a system such as

\[
\boxed{
a_n=a_{n-1}+b_{n-1}
}
\]

and

\[
\boxed{
b_n=a_{n-1}.
}
\]

Systems of linear recurrences can often be expressed using matrices.

%%%%%%%%%%%%%%%%%%%%%%%%%%%%%%%%%%%%%%%%%%%%%%%%%%%%%%%%%%%%
\subsection{Matrix Form of Recurrences}
\label{subsec:matrix-form-recurrence}
%%%%%%%%%%%%%%%%%%%%%%%%%%%%%%%%%%%%%%%%%%%%%%%%%%%%%%%%%%%%

The Fibonacci recurrence can be written

\[
\begin{pmatrix}
F_n\\
F_{n-1}
\end{pmatrix}
=
\begin{pmatrix}
1&1\\
1&0
\end{pmatrix}
\begin{pmatrix}
F_{n-1}\\
F_{n-2}
\end{pmatrix}.
\]

Thus,

\[
\boxed{
\begin{pmatrix}
F_n\\
F_{n-1}
\end{pmatrix}
=
\begin{pmatrix}
1&1\\
1&0
\end{pmatrix}^{n-1}
\begin{pmatrix}
F_1\\
F_0
\end{pmatrix}.
}
\]

This connects recurrence relations with eigenvalues, diagonalization,
and linear algebra.

%%%%%%%%%%%%%%%%%%%%%%%%%%%%%%%%%%%%%%%%%%%%%%%%%%%%%%%%%%%%
\subsection{Nonhomogeneous Linear Recurrences}
\label{subsec:nonhomogeneous-linear-recurrences}
%%%%%%%%%%%%%%%%%%%%%%%%%%%%%%%%%%%%%%%%%%%%%%%%%%%%%%%%%%%%

Consider

\[
\boxed{
a_n
=
c_1a_{n-1}
+\cdots+
c_ka_{n-k}
+
f(n).
}
\]

The general solution has the form

\[
\boxed{
a_n
=
a_n^{(h)}
+
a_n^{(p)},
}
\]

where

\[
a_n^{(h)}
\]

solves the associated homogeneous recurrence and

\[
a_n^{(p)}
\]

is one particular solution.

%%%%%%%%%%%%%%%%%%%%%%%%%%%%%%%%%%%%%%%%%%%%%%%%%%%%%%%%%%%%
\subsection{Method of Undetermined Coefficients}
\label{subsec:undetermined-coefficients}
%%%%%%%%%%%%%%%%%%%%%%%%%%%%%%%%%%%%%%%%%%%%%%%%%%%%%%%%%%%%

For certain forcing terms \(f(n)\), one can guess a particular solution
of a similar form.

Typical choices include:

\[
\begin{array}{c|c}
f(n) & \text{Typical Guess}\\
\hline
c & A\\
n & An+B\\
n^2 & An^2+Bn+C\\
\lambda^n & A\lambda^n\\
n\lambda^n & (An+B)\lambda^n
\end{array}
\]

If the proposed form duplicates a homogeneous solution, it must be
multiplied by a sufficient power of \(n\).

%%%%%%%%%%%%%%%%%%%%%%%%%%%%%%%%%%%%%%%%%%%%%%%%%%%%%%%%%%%%
\subsection{Worked Example: Constant Forcing}
\label{subsec:worked-constant-forcing}
%%%%%%%%%%%%%%%%%%%%%%%%%%%%%%%%%%%%%%%%%%%%%%%%%%%%%%%%%%%%

\begin{workedexamplebox}{Solving a Nonhomogeneous Recurrence}

Consider

\[
a_n=3a_{n-1}+4.
\]

The associated homogeneous recurrence is

\[
a_n^{(h)}=3a_{n-1}^{(h)},
\]

so

\[
a_n^{(h)}=C3^n.
\]

For a particular solution, try a constant

\[
a_n^{(p)}=A.
\]

Then

\[
A=3A+4,
\]

so

\[
A=-2.
\]

Therefore,

\[
a_n=C3^n-2.
\]

If

\[
a_0=5,
\]

then

\[
5=C-2,
\]

so

\[
C=7.
\]

\begin{answerbox}
\[
\boxed{
a_n=7\cdot3^n-2.
}
\]
\end{answerbox}

\end{workedexamplebox}

%%%%%%%%%%%%%%%%%%%%%%%%%%%%%%%%%%%%%%%%%%%%%%%%%%%%%%%%%%%%
\subsection{Resonance}
\label{subsec:recurrence-resonance}
%%%%%%%%%%%%%%%%%%%%%%%%%%%%%%%%%%%%%%%%%%%%%%%%%%%%%%%%%%%%

Suppose

\[
a_n-2a_{n-1}=2^n.
\]

The homogeneous solution contains

\[
C2^n.
\]

Therefore a particular solution of the form

\[
A2^n
\]

would duplicate the homogeneous term and fail.

Instead, try

\[
\boxed{
An2^n.
}
\]

This phenomenon is analogous to resonance in differential equations.

%%%%%%%%%%%%%%%%%%%%%%%%%%%%%%%%%%%%%%%%%%%%%%%%%%%%%%%%%%%%
\subsection{Telescoping Recurrences}
\label{subsec:telescoping-recurrences}
%%%%%%%%%%%%%%%%%%%%%%%%%%%%%%%%%%%%%%%%%%%%%%%%%%%%%%%%%%%%

Some recurrences can be rewritten so that successive differences are
easy to sum.

Suppose

\[
a_n-a_{n-1}=n.
\]

Summing from \(1\) to \(n\),

\[
\sum_{k=1}^{n}(a_k-a_{k-1})
=
\sum_{k=1}^{n}k.
\]

The left side telescopes:

\[
a_n-a_0
=
\frac{n(n+1)}{2}.
\]

Therefore,

\[
\boxed{
a_n
=
a_0+\frac{n(n+1)}{2}.
}
\]

%%%%%%%%%%%%%%%%%%%%%%%%%%%%%%%%%%%%%%%%%%%%%%%%%%%%%%%%%%%%
\subsection{Worked Example: Telescoping}
\label{subsec:worked-telescoping}
%%%%%%%%%%%%%%%%%%%%%%%%%%%%%%%%%%%%%%%%%%%%%%%%%%%%%%%%%%%%

\begin{workedexamplebox}{A Difference Recurrence}

Suppose

\[
a_0=3
\]

and

\[
a_n=a_{n-1}+2n.
\]

Then

\[
a_n-a_{n-1}=2n.
\]

Summing from \(1\) to \(n\),

\[
a_n-a_0
=
2\sum_{k=1}^{n}k.
\]

Thus,

\[
a_n-3
=
n(n+1).
\]

\begin{answerbox}
\[
\boxed{
a_n=n(n+1)+3.
}
\]
\end{answerbox}

\end{workedexamplebox}

%%%%%%%%%%%%%%%%%%%%%%%%%%%%%%%%%%%%%%%%%%%%%%%%%%%%%%%%%%%%
\subsection{Divide-and-Conquer Recurrences}
\label{subsec:divide-and-conquer-recurrences}
%%%%%%%%%%%%%%%%%%%%%%%%%%%%%%%%%%%%%%%%%%%%%%%%%%%%%%%%%%%%

Algorithms often divide a problem into smaller subproblems.

A common recurrence has the form

\[
\boxed{
T(n)
=
aT\left(\frac{n}{b}\right)
+
f(n),
}
\]

where:

\[
a
=
\text{number of subproblems},
\]

\[
\frac{n}{b}
=
\text{size of each subproblem},
\]

and

\[
f(n)
=
\text{nonrecursive work}.
\]

These are called divide-and-conquer recurrences.

%%%%%%%%%%%%%%%%%%%%%%%%%%%%%%%%%%%%%%%%%%%%%%%%%%%%%%%%%%%%
\subsection{Example: Binary Search}
\label{subsec:binary-search-recurrence}
%%%%%%%%%%%%%%%%%%%%%%%%%%%%%%%%%%%%%%%%%%%%%%%%%%%%%%%%%%%%

Binary search repeatedly halves the search interval.

A simplified running-time recurrence is

\[
\boxed{
T(n)
=
T\left(\frac n2\right)
+
1.
}
\]

After \(k\) halvings,

\[
\frac{n}{2^k}=1.
\]

Thus,

\[
2^k=n
\]

and

\[
k=\log_2 n.
\]

Therefore,

\[
\boxed{
T(n)=\Theta(\log n).
}
\]

%%%%%%%%%%%%%%%%%%%%%%%%%%%%%%%%%%%%%%%%%%%%%%%%%%%%%%%%%%%%
\subsection{Asymptotic Notation}
\label{subsec:recurrence-asymptotic-notation}
%%%%%%%%%%%%%%%%%%%%%%%%%%%%%%%%%%%%%%%%%%%%%%%%%%%%%%%%%%%%

Algorithmic recurrence relations are frequently analyzed using
asymptotic notation.

The statement

\[
T(n)=O(f(n))
\]

means that \(T(n)\) grows no faster than a constant multiple of \(f(n)\)
for sufficiently large \(n\).

The statement

\[
T(n)=\Omega(f(n))
\]

gives an asymptotic lower bound.

The statement

\[
\boxed{
T(n)=\Theta(f(n))
}
\]

means that \(f(n)\) provides both an asymptotic upper and lower bound.

The Greek capital letter

\[
\Theta
\]

is called ``Theta.''

%%%%%%%%%%%%%%%%%%%%%%%%%%%%%%%%%%%%%%%%%%%%%%%%%%%%%%%%%%%%
\subsection{Recursion Trees}
\label{subsec:recursion-trees}
%%%%%%%%%%%%%%%%%%%%%%%%%%%%%%%%%%%%%%%%%%%%%%%%%%%%%%%%%%%%

A recursion tree displays the repeated expansion of a divide-and-conquer
recurrence.

For example,

\[
T(n)=2T(n/2)+n
\]

creates two subproblems at each level.

At level \(j\):

\[
2^j
\]

subproblems each have size

\[
\frac{n}{2^j}.
\]

The total nonrecursive work on the level is

\[
2^j\frac{n}{2^j}
=
n.
\]

There are approximately

\[
\log_2 n
\]

levels.

Hence,

\[
\boxed{
T(n)=\Theta(n\log n).
}
\]

%%%%%%%%%%%%%%%%%%%%%%%%%%%%%%%%%%%%%%%%%%%%%%%%%%%%%%%%%%%%
\subsection{The Master Theorem}
\label{subsec:master-theorem}
%%%%%%%%%%%%%%%%%%%%%%%%%%%%%%%%%%%%%%%%%%%%%%%%%%%%%%%%%%%%

The Master Theorem gives asymptotic solutions to many recurrences of
the form

\[
\boxed{
T(n)
=
aT\left(\frac nb\right)
+
f(n),
}
\]

with

\[
a\geq1
\]

and

\[
b>1.
\]

The critical comparison is between

\[
f(n)
\]

and

\[
n^{\log_b a}.
\]

%%%%%%%%%%%%%%%%%%%%%%%%%%%%%%%%%%%%%%%%%%%%%%%%%%%%%%%%%%%%
\subsection{Master Theorem: Case 1}
\label{subsec:master-theorem-case-one}
%%%%%%%%%%%%%%%%%%%%%%%%%%%%%%%%%%%%%%%%%%%%%%%%%%%%%%%%%%%%

If

\[
f(n)
=
O\left(
n^{\log_b a-\varepsilon}
\right)
\]

for some

\[
\varepsilon>0,
\]

then the recursive subproblems dominate.

Therefore,

\[
\boxed{
T(n)
=
\Theta\left(
n^{\log_ba}
\right).
}
\]

The Greek letter

\[
\varepsilon
\]

is called ``epsilon.''

%%%%%%%%%%%%%%%%%%%%%%%%%%%%%%%%%%%%%%%%%%%%%%%%%%%%%%%%%%%%
\subsection{Master Theorem: Case 2}
\label{subsec:master-theorem-case-two}
%%%%%%%%%%%%%%%%%%%%%%%%%%%%%%%%%%%%%%%%%%%%%%%%%%%%%%%%%%%%

In the basic form, if

\[
f(n)
=
\Theta\left(
n^{\log_b a}\log^k n
\right)
\]

for some

\[
k\geq0,
\]

then

\[
\boxed{
T(n)
=
\Theta\left(
n^{\log_b a}
\log^{k+1}n
\right).
}
\]

The common special case

\[
f(n)
=
\Theta\left(
n^{\log_ba}
\right)
\]

gives

\[
\boxed{
T(n)
=
\Theta\left(
n^{\log_ba}\log n
\right).
}
\]

%%%%%%%%%%%%%%%%%%%%%%%%%%%%%%%%%%%%%%%%%%%%%%%%%%%%%%%%%%%%
\subsection{Master Theorem: Case 3}
\label{subsec:master-theorem-case-three}
%%%%%%%%%%%%%%%%%%%%%%%%%%%%%%%%%%%%%%%%%%%%%%%%%%%%%%%%%%%%

If

\[
f(n)
=
\Omega\left(
n^{\log_ba+\varepsilon}
\right)
\]

for some

\[
\varepsilon>0,
\]

and an appropriate regularity condition holds, then the nonrecursive
work dominates.

Therefore,

\[
\boxed{
T(n)=\Theta(f(n)).
}
\]

A common regularity condition is that there exists a constant

\[
c<1
\]

such that

\[
af(n/b)\leq cf(n)
\]

for sufficiently large \(n\).

%%%%%%%%%%%%%%%%%%%%%%%%%%%%%%%%%%%%%%%%%%%%%%%%%%%%%%%%%%%%
\subsection{Worked Example: Master Theorem}
\label{subsec:worked-master-theorem}
%%%%%%%%%%%%%%%%%%%%%%%%%%%%%%%%%%%%%%%%%%%%%%%%%%%%%%%%%%%%

\begin{workedexamplebox}{Merge-Sort-Type Recurrence}

Consider

\[
T(n)=2T(n/2)+n.
\]

Here,

\[
a=2,
\qquad
b=2,
\qquad
f(n)=n.
\]

Compute

\[
n^{\log_ba}
=
n^{\log_2 2}
=
n.
\]

Thus,

\[
f(n)
=
\Theta\left(
n^{\log_ba}
\right).
\]

This is the basic balanced case.

\begin{answerbox}
\[
\boxed{
T(n)=\Theta(n\log n).
}
\]
\end{answerbox}

\end{workedexamplebox}

%%%%%%%%%%%%%%%%%%%%%%%%%%%%%%%%%%%%%%%%%%%%%%%%%%%%%%%%%%%%
\subsection{Worked Example: Dominant Recursive Work}
\label{subsec:worked-master-case-one}
%%%%%%%%%%%%%%%%%%%%%%%%%%%%%%%%%%%%%%%%%%%%%%%%%%%%%%%%%%%%

\begin{workedexamplebox}{Applying Master Theorem Case 1}

Consider

\[
T(n)=4T(n/2)+n.
\]

We have

\[
a=4,
\qquad
b=2.
\]

Therefore,

\[
n^{\log_2 4}
=
n^2.
\]

Since

\[
n
=
O(n^{2-1}),
\]

the recursive work dominates.

\begin{answerbox}
\[
\boxed{
T(n)=\Theta(n^2).
}
\]
\end{answerbox}

\end{workedexamplebox}

%%%%%%%%%%%%%%%%%%%%%%%%%%%%%%%%%%%%%%%%%%%%%%%%%%%%%%%%%%%%
\subsection{Worked Example: Dominant Nonrecursive Work}
\label{subsec:worked-master-case-three}
%%%%%%%%%%%%%%%%%%%%%%%%%%%%%%%%%%%%%%%%%%%%%%%%%%%%%%%%%%%%

\begin{workedexamplebox}{Applying Master Theorem Case 3}

Consider

\[
T(n)=2T(n/2)+n^2.
\]

Here,

\[
n^{\log_2 2}=n.
\]

The forcing term

\[
n^2
\]

grows polynomially faster than \(n\).

Also,

\[
2\left(\frac n2\right)^2
=
\frac12n^2,
\]

so the regularity condition holds.

\begin{answerbox}
\[
\boxed{
T(n)=\Theta(n^2).
}
\]
\end{answerbox}

\end{workedexamplebox}

%%%%%%%%%%%%%%%%%%%%%%%%%%%%%%%%%%%%%%%%%%%%%%%%%%%%%%%%%%%%
\subsection{Limits of the Master Theorem}
\label{subsec:limits-master-theorem}
%%%%%%%%%%%%%%%%%%%%%%%%%%%%%%%%%%%%%%%%%%%%%%%%%%%%%%%%%%%%

The Master Theorem does not apply to every recurrence.

For example,

\[
T(n)
=
T(n-1)+n
\]

does not have the required divide-and-conquer form.

Neither does

\[
T(n)
=
T(n/2)+T(n/3)+n.
\]

Other methods include:

\begin{itemize}
    \item direct iteration;
    \item substitution;
    \item recursion trees;
    \item induction;
    \item generating functions;
    \item more general divide-and-conquer theorems.
\end{itemize}

%%%%%%%%%%%%%%%%%%%%%%%%%%%%%%%%%%%%%%%%%%%%%%%%%%%%%%%%%%%%
\subsection{Substitution Method}
\label{subsec:substitution-method-recurrence}
%%%%%%%%%%%%%%%%%%%%%%%%%%%%%%%%%%%%%%%%%%%%%%%%%%%%%%%%%%%%

The substitution method for asymptotic recurrence analysis consists of:

\begin{enumerate}
    \item guessing an asymptotic bound;
    \item assuming the bound for smaller inputs;
    \item substituting it into the recurrence;
    \item proving the bound inductively.
\end{enumerate}

For example, one might conjecture

\[
T(n)=O(n\log n)
\]

and then verify this using induction.

This approach is especially useful when the Master Theorem is
inapplicable.

%%%%%%%%%%%%%%%%%%%%%%%%%%%%%%%%%%%%%%%%%%%%%%%%%%%%%%%%%%%%
\subsection{Recurrence Verification}
\label{subsec:recurrence-verification}
%%%%%%%%%%%%%%%%%%%%%%%%%%%%%%%%%%%%%%%%%%%%%%%%%%%%%%%%%%%%

A proposed explicit formula should satisfy:

\begin{enumerate}
    \item the recurrence relation;
    \item every required initial condition.
\end{enumerate}

Both are essential.

A formula may satisfy the recurrence but correspond to a different
sequence because its initial conditions differ.

%%%%%%%%%%%%%%%%%%%%%%%%%%%%%%%%%%%%%%%%%%%%%%%%%%%%%%%%%%%%
\subsection{Worked Example: Verifying a Solution}
\label{subsec:worked-verify-recurrence}
%%%%%%%%%%%%%%%%%%%%%%%%%%%%%%%%%%%%%%%%%%%%%%%%%%%%%%%%%%%%

\begin{workedexamplebox}{Checking an Explicit Formula}

Consider

\[
a_n=2a_{n-1}+1,
\qquad
a_0=0.
\]

Suppose

\[
a_n=2^n-1.
\]

Then

\[
2a_{n-1}+1
=
2(2^{n-1}-1)+1
=
2^n-1.
\]

Thus the recurrence is satisfied.

Also,

\[
a_0
=
2^0-1
=
0.
\]

\begin{answerbox}
The formula

\[
\boxed{
a_n=2^n-1
}
\]

satisfies both the recurrence and initial condition.
\end{answerbox}

\end{workedexamplebox}

%%%%%%%%%%%%%%%%%%%%%%%%%%%%%%%%%%%%%%%%%%%%%%%%%%%%%%%%%%%%
\subsection{Recurrences and Mathematical Induction}
\label{subsec:recurrences-induction}
%%%%%%%%%%%%%%%%%%%%%%%%%%%%%%%%%%%%%%%%%%%%%%%%%%%%%%%%%%%%

Recurrence relations and induction are closely related.

A recurrence builds an object or sequence from smaller cases.

Induction proves that a property remains valid as the size increases.

Once an explicit formula is guessed, induction often provides the
simplest verification.

Thus,

\[
\boxed{
\text{recursion constructs}
\qquad
\text{while}
\qquad
\text{induction verifies}.
}
\]

%%%%%%%%%%%%%%%%%%%%%%%%%%%%%%%%%%%%%%%%%%%%%%%%%%%%%%%%%%%%
\subsection{Convolution Recurrences}
\label{subsec:convolution-recurrences}
%%%%%%%%%%%%%%%%%%%%%%%%%%%%%%%%%%%%%%%%%%%%%%%%%%%%%%%%%%%%

Some combinatorial recurrences involve sums of products.

A general form is

\[
\boxed{
a_n
=
\sum_{k=0}^{n-1}
a_k a_{n-1-k}.
}
\]

Such expressions are called convolutions.

They arise when an object is decomposed into two independent smaller
objects whose sizes add to \(n-1\).

The most famous example leads to the Catalan numbers.

%%%%%%%%%%%%%%%%%%%%%%%%%%%%%%%%%%%%%%%%%%%%%%%%%%%%%%%%%%%%
\subsection{Catalan Recurrence}
\label{subsec:catalan-recurrence}
%%%%%%%%%%%%%%%%%%%%%%%%%%%%%%%%%%%%%%%%%%%%%%%%%%%%%%%%%%%%

The Catalan numbers satisfy

\[
C_0=1
\]

and

\[
\boxed{
C_{n+1}
=
\sum_{k=0}^{n}
C_kC_{n-k}.
}
\]

This recurrence appears in many structures, including:

\begin{itemize}
    \item binary trees;
    \item balanced parenthesis strings;
    \item triangulations;
    \item certain lattice paths;
    \item polygon dissections.
\end{itemize}

The same recurrence arising from apparently different structures hints
at hidden combinatorial relationships among those structures.

%%%%%%%%%%%%%%%%%%%%%%%%%%%%%%%%%%%%%%%%%%%%%%%%%%%%%%%%%%%%
\subsection{Catalan Closed Form}
\label{subsec:catalan-closed-form-preview}
%%%%%%%%%%%%%%%%%%%%%%%%%%%%%%%%%%%%%%%%%%%%%%%%%%%%%%%%%%%%

The Catalan numbers also have the explicit form

\[
\boxed{
C_n
=
\frac{1}{n+1}
\binom{2n}{n}.
}
\]

The first values are

\[
1,1,2,5,14,42,132,\ldots.
\]

Their systematic derivation is more naturally handled using generating
functions and will be revisited later in this chapter.

%%%%%%%%%%%%%%%%%%%%%%%%%%%%%%%%%%%%%%%%%%%%%%%%%%%%%%%%%%%%
\subsection{Dynamic Programming and Recurrences}
\label{subsec:dynamic-programming-recurrences}
%%%%%%%%%%%%%%%%%%%%%%%%%%%%%%%%%%%%%%%%%%%%%%%%%%%%%%%%%%%%

Dynamic programming converts an optimization or counting problem into
recurrences among smaller subproblems.

A typical pattern is

\[
\boxed{
F(i,j)
=
\min
\left\{
F(i-1,j),
F(i,j-1),
F(i-1,j-1)+c_{ij}
\right\}.
}
\]

The exact recurrence depends on the application.

Examples include:

\begin{itemize}
    \item shortest paths;
    \item sequence alignment;
    \item knapsack;
    \item edit distance;
    \item matrix-chain multiplication;
    \item optimal scheduling.
\end{itemize}

The recurrence supplies the mathematical structure, while dynamic
programming provides an efficient computational strategy.

%%%%%%%%%%%%%%%%%%%%%%%%%%%%%%%%%%%%%%%%%%%%%%%%%%%%%%%%%%%%
\subsection{Memoization}
\label{subsec:memoization}
%%%%%%%%%%%%%%%%%%%%%%%%%%%%%%%%%%%%%%%%%%%%%%%%%%%%%%%%%%%%

A naive recursive program may repeatedly solve the same subproblem.

Memoization stores previously computed results.

Thus if the recurrence requests the same value again, the stored result
is reused rather than recomputed.

This may reduce exponential recursive computation to polynomial time.

%%%%%%%%%%%%%%%%%%%%%%%%%%%%%%%%%%%%%%%%%%%%%%%%%%%%%%%%%%%%
\subsection{Bottom-Up Computation}
\label{subsec:bottom-up-recurrence}
%%%%%%%%%%%%%%%%%%%%%%%%%%%%%%%%%%%%%%%%%%%%%%%%%%%%%%%%%%%%

Instead of recursively requesting earlier values, one may compute the
sequence from smallest to largest indices.

For Fibonacci numbers:

\[
F_0=0,
\qquad
F_1=1,
\]

then compute

\[
F_2,
F_3,
\ldots,
F_n.
\]

This is a bottom-up dynamic programming approach.

It often avoids recursion overhead and makes dependency structure
explicit.

%%%%%%%%%%%%%%%%%%%%%%%%%%%%%%%%%%%%%%%%%%%%%%%%%%%%%%%%%%%%
\subsection{State Design}
\label{subsec:recurrence-state-design}
%%%%%%%%%%%%%%%%%%%%%%%%%%%%%%%%%%%%%%%%%%%%%%%%%%%%%%%%%%%%

In combinatorial recurrence modeling, one of the hardest tasks is
choosing the correct state.

A state should contain enough information to determine which extensions
are allowed.

For example, counting binary strings with no consecutive \(1\)'s may
require tracking whether the current string ends in \(0\) or \(1\).

If too little information is stored, the recurrence is invalid.

If too much information is stored, the recurrence may become
unnecessarily complicated.

%%%%%%%%%%%%%%%%%%%%%%%%%%%%%%%%%%%%%%%%%%%%%%%%%%%%%%%%%%%%
\subsection{Boundary Conditions}
\label{subsec:boundary-conditions-recurrence}
%%%%%%%%%%%%%%%%%%%%%%%%%%%%%%%%%%%%%%%%%%%%%%%%%%%%%%%%%%%%

Combinatorial recurrences frequently require carefully chosen base
cases.

For example, when counting tilings,

\[
T_0=1
\]

may appear unusual at first.

However, there is exactly one way to tile an empty board:

\[
\boxed{
\text{use no tiles}.
}
\]

This convention allows recurrence formulas to remain valid without
unnecessary exceptions.

%%%%%%%%%%%%%%%%%%%%%%%%%%%%%%%%%%%%%%%%%%%%%%%%%%%%%%%%%%%%
\subsection{Recurrences and Generating Functions}
\label{subsec:recurrences-generating-functions-preview}
%%%%%%%%%%%%%%%%%%%%%%%%%%%%%%%%%%%%%%%%%%%%%%%%%%%%%%%%%%%%

A recurrence can often be converted into an algebraic equation involving
a generating function.

For a sequence

\[
a_0,a_1,a_2,\ldots,
\]

define

\[
A(x)
=
\sum_{n=0}^{\infty}a_nx^n.
\]

Then shifts such as

\[
a_{n-1}
\]

correspond to multiplication by \(x\).

This allows recurrence equations to become algebraic equations for
\(A(x)\).

Generating functions will be developed systematically in the next
section.

%%%%%%%%%%%%%%%%%%%%%%%%%%%%%%%%%%%%%%%%%%%%%%%%%%%%%%%%%%%%
\subsection{Growth from Characteristic Roots}
\label{subsec:growth-characteristic-roots}
%%%%%%%%%%%%%%%%%%%%%%%%%%%%%%%%%%%%%%%%%%%%%%%%%%%%%%%%%%%%

For a homogeneous constant-coefficient recurrence, the largest
characteristic root in magnitude often controls long-term growth.

Suppose

\[
a_n
=
A r_1^n
+
B r_2^n
\]

and

\[
|r_1|>|r_2|.
\]

If \(A\neq0\), then asymptotically

\[
\boxed{
a_n
\sim
A r_1^n.
}
\]

The root of largest magnitude is called a dominant root.

%%%%%%%%%%%%%%%%%%%%%%%%%%%%%%%%%%%%%%%%%%%%%%%%%%%%%%%%%%%%
\subsection{Oscillatory Behavior}
\label{subsec:oscillatory-recurrences}
%%%%%%%%%%%%%%%%%%%%%%%%%%%%%%%%%%%%%%%%%%%%%%%%%%%%%%%%%%%%

Negative or complex characteristic roots can produce oscillation.

For example,

\[
a_n=-a_{n-1}
\]

with

\[
a_0=1
\]

gives

\[
1,-1,1,-1,\ldots.
\]

The explicit formula is

\[
\boxed{
a_n=(-1)^n.
}
\]

Complex roots similarly produce sine and cosine behavior.

%%%%%%%%%%%%%%%%%%%%%%%%%%%%%%%%%%%%%%%%%%%%%%%%%%%%%%%%%%%%
\subsection{Stable and Unstable First-Order Recurrences}
\label{subsec:stable-unstable-recurrences}
%%%%%%%%%%%%%%%%%%%%%%%%%%%%%%%%%%%%%%%%%%%%%%%%%%%%%%%%%%%%

Consider

\[
a_n=ra_{n-1}+b.
\]

The equilibrium is

\[
L=\frac{b}{1-r}.
\]

Since

\[
a_n-L
=
r^n(a_0-L),
\]

the long-term behavior depends on \(|r|\).

If

\[
|r|<1,
\]

then

\[
a_n\to L.
\]

If

\[
|r|>1,
\]

deviations from \(L\) generally grow.

If

\[
r=-1,
\]

oscillation may occur.

This provides an elementary connection to discrete dynamical systems.

%%%%%%%%%%%%%%%%%%%%%%%%%%%%%%%%%%%%%%%%%%%%%%%%%%%%%%%%%%%%
\subsection{Recurrence Trees in Combinatorics}
\label{subsec:recurrence-trees-combinatorics}
%%%%%%%%%%%%%%%%%%%%%%%%%%%%%%%%%%%%%%%%%%%%%%%%%%%%%%%%%%%%

Recursion trees are not limited to algorithm analysis.

A recursive combinatorial class can often be represented by a tree in
which:

\begin{itemize}
    \item the root is the original object;
    \item children correspond to possible decompositions;
    \item leaves correspond to base cases.
\end{itemize}

Counting leaves or weights across levels can reveal recurrence
solutions.

%%%%%%%%%%%%%%%%%%%%%%%%%%%%%%%%%%%%%%%%%%%%%%%%%%%%%%%%%%%%
\subsection{Common Errors}
\label{subsec:recurrence-common-errors}
%%%%%%%%%%%%%%%%%%%%%%%%%%%%%%%%%%%%%%%%%%%%%%%%%%%%%%%%%%%%

\subsubsection{Forgetting Initial Conditions}

A recurrence relation usually does not define a unique sequence without
appropriate starting values.

%%%%%%%%%%%%%%%%%%%%%%%%%%%%%%%%%%%%%%%%%%%%%%%%%%%%%%%%%%%%
\subsubsection{Using Too Few Initial Conditions}

A second-order recurrence typically requires two starting values.

A \(k\)-th order recurrence generally requires \(k\) independent initial
conditions.

%%%%%%%%%%%%%%%%%%%%%%%%%%%%%%%%%%%%%%%%%%%%%%%%%%%%%%%%%%%%
\subsubsection{Using the Wrong Characteristic Polynomial}

For

\[
a_n
=
c_1a_{n-1}
+
c_2a_{n-2},
\]

the characteristic equation is

\[
\boxed{
r^2-c_1r-c_2=0.
}
\]

The signs must be handled carefully.

%%%%%%%%%%%%%%%%%%%%%%%%%%%%%%%%%%%%%%%%%%%%%%%%%%%%%%%%%%%%
\subsubsection{Ignoring Repeated Roots}

A repeated root does not produce two copies of the same term.

For a repeated root \(r\), the second independent solution is

\[
nr^n.
\]

%%%%%%%%%%%%%%%%%%%%%%%%%%%%%%%%%%%%%%%%%%%%%%%%%%%%%%%%%%%%
\subsubsection{Forgetting Resonance}

If a proposed particular solution duplicates a homogeneous solution,
multiply the trial form by a sufficient power of \(n\).

%%%%%%%%%%%%%%%%%%%%%%%%%%%%%%%%%%%%%%%%%%%%%%%%%%%%%%%%%%%%
\subsubsection{Solving the Recurrence but Ignoring Base Cases}

A formula that satisfies the recurrence may still fail the required
initial conditions.

Always verify both.

%%%%%%%%%%%%%%%%%%%%%%%%%%%%%%%%%%%%%%%%%%%%%%%%%%%%%%%%%%%%
\subsubsection{Applying the Master Theorem Outside Its Domain}

The Master Theorem applies only to particular divide-and-conquer
recurrences.

Do not force a recurrence such as

\[
T(n)=T(n-1)+n
\]

into the theorem.

%%%%%%%%%%%%%%%%%%%%%%%%%%%%%%%%%%%%%%%%%%%%%%%%%%%%%%%%%%%%
\subsubsection{Confusing Exact and Asymptotic Solutions}

The statement

\[
T(n)=\Theta(n\log n)
\]

describes growth rate.

It does not specify the exact value of \(T(n)\).

%%%%%%%%%%%%%%%%%%%%%%%%%%%%%%%%%%%%%%%%%%%%%%%%%%%%%%%%%%%%
\subsubsection{Incorrect Combinatorial Decomposition}

When deriving a recurrence by cases, the cases must:

\[
\boxed{
\text{cover every object}
}
\]

and usually should be disjoint.

Otherwise the recurrence may undercount or overcount.

%%%%%%%%%%%%%%%%%%%%%%%%%%%%%%%%%%%%%%%%%%%%%%%%%%%%%%%%%%%%
\subsubsection{Using an Incomplete State}

If future choices depend on information not stored in the recurrence
state, the recurrence is not valid.

%%%%%%%%%%%%%%%%%%%%%%%%%%%%%%%%%%%%%%%%%%%%%%%%%%%%%%%%%%%%
\subsection{Applications}
\label{subsec:recurrence-applications}
%%%%%%%%%%%%%%%%%%%%%%%%%%%%%%%%%%%%%%%%%%%%%%%%%%%%%%%%%%%%

\begin{applicationbox}

\textbf{Combinatorics.}
Recurrences count strings, tilings, trees, partitions, lattice paths,
and many recursively defined structures.

\medskip

\textbf{Computer Science.}
Algorithm running times, divide-and-conquer methods, recursive
procedures, and dynamic programming are modeled by recurrence
relations.

\medskip

\textbf{Probability.}
Expected values and probabilities in discrete stochastic processes
often satisfy recurrences.

\medskip

\textbf{Finance.}
Repeated deposits, loans, interest accumulation, and discrete cash-flow
models naturally produce first-order recurrences.

\medskip

\textbf{Population Models.}
Discrete-time population growth may be modeled by linear or nonlinear
recurrences.

\medskip

\textbf{Coding and Information Theory.}
Restricted strings and finite-state encodings often generate recurrence
relations.

\medskip

\textbf{Operations Research.}
Dynamic programming recurrences arise in routing, scheduling, inventory,
resource allocation, and sequential decision problems.

\medskip

\textbf{Discrete Dynamical Systems.}
Iterated maps

\[
x_{n+1}=f(x_n)
\]

are recurrence relations and form the basic language of discrete-time
dynamics.

\end{applicationbox}

%%%%%%%%%%%%%%%%%%%%%%%%%%%%%%%%%%%%%%%%%%%%%%%%%%%%%%%%%%%%
\subsection{Exercises}
\label{subsec:recurrence-exercises}
%%%%%%%%%%%%%%%%%%%%%%%%%%%%%%%%%%%%%%%%%%%%%%%%%%%%%%%%%%%%

\begin{exercise}[1]
Compute the first five terms of the sequence defined by

\[
a_0=2,
\qquad
a_n=a_{n-1}+4.
\]
\end{exercise}

\begin{exercise}[1]
Compute the first six terms of

\[
a_0=3,
\qquad
a_n=2a_{n-1}.
\]
\end{exercise}

\begin{exercise}[1]
State the order of

\[
a_n=4a_{n-1}-a_{n-3}.
\]
\end{exercise}

\begin{exercise}[1]
Determine whether

\[
a_n=2a_{n-1}+5
\]

is homogeneous or nonhomogeneous.
\end{exercise}

\begin{exercise}[1]
Determine whether

\[
a_n=(a_{n-1})^2+1
\]

is linear or nonlinear.
\end{exercise}

\begin{exercise}[1]
Solve

\[
a_n=a_{n-1}+5,
\qquad
a_0=7.
\]
\end{exercise}

\begin{exercise}[1]
Solve

\[
a_n=4a_{n-1},
\qquad
a_0=2.
\]
\end{exercise}

\begin{exercise}[1]
Verify that

\[
a_n=3\cdot2^n-1
\]

satisfies

\[
a_n=2a_{n-1}+1.
\]
\end{exercise}

\begin{exercise}[1]
Write the first eight Fibonacci numbers beginning with

\[
F_0=0,
\qquad
F_1=1.
\]
\end{exercise}

\begin{exercise}[2]
Solve

\[
a_n=3a_{n-1}+2,
\qquad
a_0=1.
\]
\end{exercise}

\begin{exercise}[2]
Solve

\[
a_n=5a_{n-1}-6a_{n-2},
\]

with

\[
a_0=2,
\qquad
a_1=7.
\]
\end{exercise}

\begin{exercise}[2]
Solve

\[
a_n=4a_{n-1}-4a_{n-2},
\]

with

\[
a_0=1,
\qquad
a_1=4.
\]
\end{exercise}

\begin{exercise}[2]
Find the characteristic equation of

\[
a_n
=
7a_{n-1}
-
10a_{n-2}.
\]
\end{exercise}

\begin{exercise}[2]
Find the characteristic roots of

\[
r^2-7r+10=0.
\]
\end{exercise}

\begin{exercise}[2]
Use the characteristic-root method to solve

\[
a_n
=
7a_{n-1}
-
10a_{n-2}
\]

given suitable initial conditions of your choice.
\end{exercise}

\begin{exercise}[2]
Solve

\[
a_n-a_{n-1}=3n,
\qquad
a_0=2.
\]
\end{exercise}

\begin{exercise}[2]
Derive a recurrence for the number of ways to tile a \(1\times n\)
board with \(1\times1\) squares and \(1\times2\) dominoes.
\end{exercise}

\begin{exercise}[2]
Derive a recurrence for binary strings of length \(n\) containing no
two consecutive \(1\)'s.
\end{exercise}

\begin{exercise}[2]
Compute the first six terms of the sequence from the previous exercise.
\end{exercise}

\begin{exercise}[2]
Find the asymptotic order of

\[
T(n)=T(n/2)+1.
\]
\end{exercise}

\begin{exercise}[2]
Use the Master Theorem to analyze

\[
T(n)=2T(n/2)+n.
\]
\end{exercise}

\begin{exercise}[2]
Use the Master Theorem to analyze

\[
T(n)=4T(n/2)+n.
\]
\end{exercise}

\begin{exercise}[2]
Use the Master Theorem to analyze

\[
T(n)=2T(n/2)+n^2.
\]
\end{exercise}

\begin{exercise}[2]
Explain why the Master Theorem does not directly apply to

\[
T(n)=T(n-1)+n.
\]
\end{exercise}

\begin{exercise}[3]
Prove by induction that the solution to

\[
a_n=a_{n-1}+d,
\qquad
a_0=c
\]

is

\[
a_n=c+nd.
\]
\end{exercise}

\begin{exercise}[3]
Prove by induction that the solution to

\[
a_n=ra_{n-1},
\qquad
a_0=c
\]

is

\[
a_n=cr^n.
\]
\end{exercise}

\begin{exercise}[3]
Derive Binet's formula for the Fibonacci sequence using characteristic
roots.
\end{exercise}

\begin{exercise}[3]
Show from Binet's formula that

\[
F_n
\sim
\frac{\phi^n}{\sqrt5}.
\]
\end{exercise}

\begin{exercise}[3]
Solve

\[
a_n=6a_{n-1}-9a_{n-2},
\]

with

\[
a_0=1,
\qquad
a_1=6.
\]
\end{exercise}

\begin{exercise}[3]
Find a particular solution of

\[
a_n-2a_{n-1}=3.
\]
\end{exercise}

\begin{exercise}[3]
Find a suitable trial form for a particular solution of

\[
a_n-2a_{n-1}=2^n.
\]
\end{exercise}

\begin{exercise}[3]
Solve

\[
a_n=2a_{n-1}+n,
\qquad
a_0=0.
\]
\end{exercise}

\begin{exercise}[3]
Represent the Fibonacci recurrence using a \(2\times2\) matrix.
\end{exercise}

\begin{exercise}[3]
Show that

\[
\begin{pmatrix}
F_{n+1}\\
F_n
\end{pmatrix}
=
\begin{pmatrix}
1&1\\
1&0
\end{pmatrix}^{n}
\begin{pmatrix}
1\\
0
\end{pmatrix}.
\]
\end{exercise}

\begin{exercise}[3]
Use a recursion tree to derive

\[
T(n)=\Theta(n\log n)
\]

for

\[
T(n)=2T(n/2)+n.
\]
\end{exercise}

\begin{exercise}[3]
Analyze

\[
T(n)=3T(n/2)+n
\]

using the Master Theorem.
\end{exercise}

\begin{exercise}[3]
Analyze

\[
T(n)=8T(n/2)+n^2
\]

using the Master Theorem.
\end{exercise}

\begin{exercise}[3]
Analyze

\[
T(n)=2T(n/2)+n\log n.
\]
\end{exercise}

\begin{exercise}[3]
Derive a recurrence for the number of strings over

\[
\{0,1,2\}
\]

that contain no consecutive \(2\)'s.
\end{exercise}

\begin{exercise}[3]
Derive a system of recurrences by separating strings according to their
final symbol.
\end{exercise}

\begin{exercise}[3]
Derive the Catalan recurrence by decomposing a full binary tree at its
root.
\end{exercise}

\begin{exercise}[3]
Verify directly that

\[
C_0=1,
\quad
C_1=1,
\quad
C_2=2,
\quad
C_3=5
\]

satisfy the Catalan recurrence.
\end{exercise}

\begin{exercise}[4]
Investigate the solution of linear homogeneous recurrences of arbitrary
order using characteristic polynomials.
\end{exercise}

\begin{exercise}[4]
Study complex characteristic roots and derive the real-valued
sine-cosine form of the solution.
\end{exercise}

\begin{exercise}[4]
Investigate repeated characteristic roots of multiplicity greater than
two.
\end{exercise}

\begin{exercise}[4]
Study generating-function solutions of the Fibonacci recurrence.
\end{exercise}

\begin{exercise}[4]
Use generating functions to derive the closed form

\[
C_n=\frac1{n+1}\binom{2n}{n}
\]

for the Catalan numbers.
\end{exercise}

\begin{exercise}[4]
Investigate linear recurrences using companion matrices and eigenvalues.
\end{exercise}

\begin{exercise}[4]
Study the connection between dominant characteristic roots and
asymptotic sequence growth.
\end{exercise}

\begin{exercise}[4]
Investigate the Akra--Bazzi method as a generalization of the Master
Theorem.
\end{exercise}

\begin{exercise}[4]
Study recurrence relations arising in dynamic programming for
matrix-chain multiplication.
\end{exercise}

\begin{exercise}[4]
Investigate nonlinear recurrences arising from the logistic map

\[
x_{n+1}=rx_n(1-x_n).
\]

Describe fixed points and basic stability behavior.
\end{exercise}

%%%%%%%%%%%%%%%%%%%%%%%%%%%%%%%%%%%%%%%%%%%%%%%%%%%%%%%%%%%%
\subsection{Conceptual Summary}
\label{subsec:recurrence-summary}
%%%%%%%%%%%%%%%%%%%%%%%%%%%%%%%%%%%%%%%%%%%%%%%%%%%%%%%%%%%%

A recurrence relation defines a sequence using previous values.

A first-order recurrence may take the form

\[
\boxed{
a_n=f(a_{n-1},n).
}
\]

A linear homogeneous recurrence with constant coefficients has the form

\[
\boxed{
a_n
=
c_1a_{n-1}
+\cdots+
c_ka_{n-k}.
}
\]

Seeking solutions of the form

\[
a_n=r^n
\]

produces a characteristic polynomial.

For a second-order recurrence

\[
a_n=c_1a_{n-1}+c_2a_{n-2},
\]

the characteristic equation is

\[
\boxed{
r^2-c_1r-c_2=0.
}
\]

Distinct roots \(r_1,r_2\) produce

\[
\boxed{
a_n
=
Ar_1^n+Br_2^n.
}
\]

A repeated root \(r\) produces

\[
\boxed{
a_n
=
(A+Bn)r^n.
}
\]

Nonhomogeneous recurrences have solutions of the form

\[
\boxed{
a_n
=
a_n^{(h)}
+
a_n^{(p)}.
}
\]

Combinatorial recurrences arise by decomposing structures into smaller
cases.

The Fibonacci recurrence

\[
\boxed{
F_n=F_{n-1}+F_{n-2}
}
\]

appears in strings, tilings, and many other discrete structures.

Convolution recurrences such as

\[
\boxed{
C_{n+1}
=
\sum_{k=0}^{n}
C_kC_{n-k}
}
\]

lead to the Catalan numbers.

Algorithmic divide-and-conquer recurrences often take the form

\[
\boxed{
T(n)
=
aT(n/b)+f(n).
}
\]

The Master Theorem compares

\[
f(n)
\]

with

\[
n^{\log_ba}.
\]

Recurrence relations therefore provide a common language connecting

\[
\boxed{
\text{combinatorics}
+
\text{algorithms}
+
\text{sequences}
+
\text{dynamic programming}
+
\text{discrete dynamics}.
}
\]

%%%%%%%%%%%%%%%%%%%%%%%%%%%%%%%%%%%%%%%%%%%%%%%%%%%%%%%%%%%%
\subsection{Section Connections}
\label{subsec:recurrence-connections}
%%%%%%%%%%%%%%%%%%%%%%%%%%%%%%%%%%%%%%%%%%%%%%%%%%%%%%%%%%%%

\chapterconnection{
Recurrence Relations build directly on the counting methods of
Section~\ref{sec:basic-counting}. Many combinatorial classes are counted
by decomposing objects of size \(n\) into smaller objects, producing
recurrences automatically. The next section, Generating Functions,
provides an algebraic method for solving and analyzing recurrence
relations. Enumerative Combinatorics will use recurrences for classical
families such as Catalan, Stirling, and partition-related sequences.
Graph Theory uses recurrences in trees, walks, and graph algorithms.
Probabilistic Combinatorics uses recurrence methods in expected-value
and random-process calculations. Optimization and Operations Research
use dynamic programming recurrences extensively, while Differential
Equations and Dynamical Systems develops the continuous and discrete
evolution viewpoint in greater depth.
}

%%%%%%%%%%%%%%%%%%%%%%%%%%%%%%%%%%%%%%%%%%%%%%%%%%%%%%%%%%%%
% SECTION SOURCE TRACKING
%%%%%%%%%%%%%%%%%%%%%%%%%%%%%%%%%%%%%%%%%%%%%%%%%%%%%%%%%%%%

%------------------------------------------------------------
% 6.3 Recurrence Relations
%------------------------------------------------------------
%
% Source basis:
%   - Standard discrete mathematics.
%   - Standard combinatorics.
%   - Standard linear recurrence theory.
%   - Standard introductory algorithm analysis.
%
% Used for:
%   - recurrence relations;
%   - initial conditions;
%   - recurrence order;
%   - recursive versus explicit descriptions;
%   - linear and nonlinear recurrences;
%   - homogeneous and nonhomogeneous recurrences;
%   - arithmetic and geometric recurrences;
%   - first-order linear recurrences;
%   - iteration;
%   - characteristic equations;
%   - distinct and repeated characteristic roots;
%   - higher-order recurrences;
%   - complex characteristic roots;
%   - Fibonacci numbers;
%   - Binet's formula;
%   - combinatorial recurrence derivation;
%   - restricted binary strings;
%   - tiling recurrences;
%   - systems of recurrences;
%   - matrix formulations;
%   - particular solutions;
%   - undetermined coefficients;
%   - resonance;
%   - telescoping recurrences;
%   - divide-and-conquer recurrences;
%   - asymptotic notation;
%   - recursion trees;
%   - Master Theorem;
%   - substitution method;
%   - solution verification;
%   - induction;
%   - convolution recurrences;
%   - Catalan recurrence;
%   - dynamic programming;
%   - memoization;
%   - bottom-up computation;
%   - recurrence-state design;
%   - generating-function preview;
%   - dominant-root growth;
%   - discrete dynamical behavior;
%   - applications and exercises.
%
% Notes:
%   - Generating-function methods are developed in Section 6.4.
%   - Catalan and other classical sequences are developed more fully in
%     Enumerative Combinatorics.
%   - Algorithmic recurrence analysis is introduced here only to the
%     level needed for discrete mathematics; computational complexity
%     and algorithms may be developed further elsewhere in the text.
%   - Matrix methods connect this section with Chapter 1 Algebra.
%
%%%%%%%%%%%%%%%%%%%%%%%%%%%%%%%%%%%%%%%%%%%%%%%%%%%%%%%%%%%%